\section{Entwicklung der Signalverarbeitungsbibliothek}

\subsection{Anwendung als DSV}

Die DSV-Bibliothek wird als eine Klasse namens Signal implementiert.
Die Klasse ist eine Fluent API\cite{fluentAPI}. Das bedeutet, dass die
nicht abschließenden Funktionen eine Referenz auf das Signalobjekt zurückgeben. 
Dadurch können Methodenaufrufe verkettet werden, um komplexe Operationen auf 
das Signal anzuwenden.

Im Folgenden wird die Entwicklung der Signalklasse Schritt für Schritt
beschrieben. Bei jedem Schritt wird der angegebene Codeblock erweitert,
bis am Ende die Klasse vorliegt, wie sie zur Zeit der Arbeit vorliegt.

\subsubsection{Die Klasse}
\label{chap:signalklasse}

Ein digitales Signal wird in seiner simpelsten Form als eine sich ständig
ändernde Variable dargestellt. Der Datentyp der Variable wird dem Nutzer 
überlassen, indem die Signalklasse als Template implementiert wird. 

\begin{lstlisting}[language=C++, 
                       frame=single, 
                       caption={Signalklasse Version 0.1},
                       captionpos=b]
template<typename T>
class Signal {
  T value;
}
\end{lstlisting}

Da Funktionen, wie das Mitteln über die letzten Werte, auf die vergangenen 
Werte zugreifen, lohnt es sich, diese ebenfalls zu speichern.
Dafür bietet sich ein RingBuffer (circular buffer) \cite{ringbuffer} an. 
Ein Ringbuffer ist ein fest definierter
Speicherblock (zum Beispiel ein Array) mit fester Größe. Das ist vor allem
auf Mikrocontrollern sehr 
wichtig, da hier Speicher und Rechenkapazität begrenzt sind. Ein Ringbuffer hat außerdem
einen Index Wert welcher anzeigt, wo in diesen Speicherblock geschrieben 
werden soll. Der Index wird bei jedem schreiben erhöht und fängt wieder am Anfang
des Speicherblocks an zu schreiben, wenn das Ende erreicht ist.
Die Größe des Speichers soll auch den Nutzer überlassen werden.

\begin{lstlisting}[language=C++, 
                       frame=single, 
                       caption={Signalklasse Version 0.2},
                       captionpos=b]
template <typename T, unsigned int BUFFER_SIZE = 256>
class Signal {
private:
  T buffer[BUFFER_SIZE] = {0};
  unsigned int index = 0;
}
\end{lstlisting}

Wenn das Signal zu einem neuen Wert übergeht, muss der Index erhöht werden.
Das passiert in einer Inputfunktion, welche mit dem Prefix \glqq from-\grqq\ 
gekennzeichnet sind. Eine Inputfunktion gibt immer eine Referenz auf das 
Signalobjekt zurück.

\begin{lstlisting}[language=C++, 
                       frame=single, 
                       caption={Signalklasse Version 0.3},
                       captionpos=b]
template <typename T, unsigned int BUFFER_SIZE = 256>
class Signal {
private:
  T buffer[BUFFER_SIZE] = {0};
  unsigned int index = 0;
public:
  Signal& fromValue(T value) {
    index = (index + 1) % BUFFER_SIZE;
    buffer[index] = value;
    return *this;
  }
}
\end{lstlisting}

Eine Funktion, die mit dem Prefix \glqq return-\grqq\ 
gekennzeichnet sind, ist eine Ausgangsfunktion. 
Solche Funktionen geben keine Signalreferenz zurück, sondern zum Beispiel den 
momentanen Wert.

\begin{lstlisting}[language=C++, 
                       frame=single, 
                       caption={Signalklasse Version 0.4},
                       captionpos=b]
template <typename T, unsigned int BUFFER_SIZE = 256>
class Signal {
private:
  T buffer[BUFFER_SIZE] = {0};
  unsigned int index = 0;
public:
  ...

  T returnValue() {
    return buffer[index];
  }
}
\end{lstlisting}

Alle anderen Funktionen sind Funktionen, welche eine Signalreferenz zurückgeben
und im Gegensatz zu Inputfunktionen nicht den Index erhöhen. Das sind zum Beispiel
add, multiply oder applyFunction. Letzteres erlaubt es dem Nutzer selbstdefinierte 
Funktionen in die Verarbeitungskette einzufügen.

\begin{lstlisting}[language=C++, 
                       frame=single, 
                       caption={Signalklasse Version 0.5},
                       captionpos=b]
template <typename T, unsigned int BUFFER_SIZE = 256>
class Signal {
private:
  T buffer[BUFFER_SIZE] = {0};
  unsigned int index = 0;
public:
  ...

  Signal& add(Signal &other) {
    this->buffer[index] += other.buffer[other.index];
    return *this;
  }

  Signal& multiply(T value) {
    this->buffer[index] = T(this->buffer[index] * value);
    return *this;
  }

  SIgnal applyFunction(T (*func)(T data, void *param), 
                       void *param) {

    buffer[index] = func(buffer[index], param);
    return *this;
  }
}
\end{lstlisting}


\subsubsection{Anwendung}
% wohin soll das Example?
Im Anhang sieht man eine Beispielanwendung in der ArduinoIDE. Hier werden ein
lichtabhängiger Widerstand und ein Potentiometer eingelesen. Der Potentiometerwert
wird dazu verwendet, um den aktuell stabilen Lichtwert zu kompensieren. Dann wird
eine LED angeschaltet, wenn der kompensierte Lichtwert unter 0 sinkt. Hält man den
lichtabhängigen Widerstand gegen eine Lichtquelle oder einen reflektierenden
Hintergrund kann durch Blocken des Lichts eine Bewegung erkannt werden. 
(Es gibt bessere Möglichkeiten dies umzusetzen aber zu Demonstrationszwecken 
sollte dieses Beispiel reichen).

Der Code verkürzt sich hier nicht wirklich. Im Gegenteil er wird wegen der 
zusätzlich definierten Funktionen erstmal komplexer. Sollen nun weitere 
Funktionen wie zum Beispiel Filter in die Verarbeitungskette eingebracht werden,
kann dies mit der Signalbibliothek einfacher und schöner umgesetzt werden.

Ein Besseres und ausfürlicheres Beispiel liefert die Erweiterung von Signal zu 
Sound im nächsten Unterkapitel und die Anwendung der Soundklasse als Synthesizer
im nächsetn Kapitel.

\subsubsection{Anmerkungen}
Ein großer Unterschied zwischen der Signalklasse und Java-streams ist, dass 
die Inputfunktionen der Signalklasse keine statischen Memberfunktionen sind.
Das liegt daran, dass ein Signal im Normalfall in einer Endlosschleife 
aktualisiert wird und nicht wie ein Stream einmal erschaffen und direkt 
beendet wird. 

Die Abtastrate hängt davon ab, wie oft pro Sekunde eine Inputfunktion
augerufen wird.

Ein weiterer wichtiger Punkt ist, dass der Index in den Inputfunktionen erhöht
werden muss, da diese der auch einen neuen Wert setzt und für andere Funktionen 
nicht garantiert ist, dass diese aufgerufen werden.

Außerdem ist es wichtig, dass nur Referenzen des Signals zurück- und übergeben werden. 
Auch wenn das Signal nicht verändert wird ist dies wichtig, da so zeitintensieves
Kopieren des Buffers verhindert wird.

